\section{Comparison of Approaches}
We compare collaborative file synchronization tools by setup effort, server requirements, and capability to synchronize heterogeneous files across local meshes \cite{googledrive,dropbox,gitbook,github,syncthing,resilio}.


\subsection{Comparison Dimensions}

The approaches introduced in the previous section address
different aspects of collaborative file management. To compare
their suitability for the problem considered in this project,
we evaluate them according to five key dimensions:

\begin{enumerate}
    \item \textbf{Membership Management.}
    The ability to discover peers, support dynamic participation, and establish collaborative groups without extensive manual configuration or centralized management.

    \item \textbf{Synchronization Model.}
    The ability to propagate changes automatically across multiple devices, support asynchronous collaboration, tolerate temporary disconnections, and maintain shared directory consistency.

    \item \textbf{History Management.}
    The ability to preserve previous file states, provide rollback mechanisms, and manage conflicting modifications made by different collaborators.

    \item \textbf{File Support.}
    The ability to handle heterogeneous file types, including source code, documents, images, multimedia content, datasets, and other binary files, as well as support efficient transfer and synchronization of large files.

    \item \textbf{Filesystem Transparency and Usability.}
    The extent to which users can interact with shared files through ordinary filesystem operations without requiring specialized workflows, commands, or collaboration tools.
\end{enumerate}

These dimensions capture the core requirements identified in
the problem definition and provide a common basis for comparing
existing approaches.


\subsection{Centralized Collaboration Platforms}

Representative examples of centralized collaboration platforms
include \textit{Google Drive}, \textit{Microsoft OneDrive},
\textit{Dropbox}, and \textit{Nextcloud}. \textit{Google Drive}
and \textit{OneDrive} are cloud services operated by large
providers and offer tightly integrated ecosystems for document
sharing and synchronization. \textit{Dropbox} focuses primarily
on cross-device file synchronization and collaborative
storage, while \textit{Nextcloud} provides a self-hosted
alternative that allows organizations and users to deploy and
manage their own collaboration infrastructure.

\subsubsection{Membership Management}
Centralized collaboration platforms generally provide strong
membership management through user accounts, invitations,
permission settings, and access control mechanisms.
Collaborative groups can be created and managed easily, but
participation is typically governed by a central authority or
service provider.

For example, \textit{Google Drive} and \textit{OneDrive}
allow users to invite collaborators through email-based
accounts and assign different access permissions.

\subsubsection{Synchronization Model}
These platforms provide automatic synchronization across
multiple devices by maintaining a central copy of shared data.
Updates are propagated through the cloud service and become
available to collaborators once synchronization is completed.
Such systems support multi-device collaboration, although they
generally depend on connectivity to the central service.

For example, a document modified in \textit{Google Drive}
is uploaded to the cloud service and subsequently synchronized
to other authorized devices connected to the same workspace.

\subsubsection{History Management}
Most centralized collaboration platforms provide basic version
history and file recovery capabilities. Previous revisions can
often be restored, but history management is usually less
comprehensive than that offered by dedicated version control
systems.

For example, \textit{Google Drive} maintains previous
revisions of documents and allows users to restore earlier
versions when necessary.

\subsubsection{File Support}
Centralized collaboration platforms support a wide range of
file types, including source code, documents, images,
multimedia content, datasets, and other binary assets. They
also support large files, although storage quotas and service
limitations may affect practical usage.

For example, \textit{Dropbox} and \textit{Nextcloud}
can store source code, documents, images, videos, and other
binary assets within the same shared workspace.

\subsubsection{Filesystem Transparency and Usability}
These systems generally provide strong filesystem transparency.
Users interact with shared files through familiar folder-based
interfaces and synchronization clients, allowing collaborative
workflows to be integrated into normal file management
practices with minimal additional training.

For example, \textit{Dropbox} provides a synchronized local
folder that can be accessed through standard filesystem
operations without requiring users to interact directly with
the underlying synchronization mechanisms.

In summary, centralized collaboration platforms provide
strong support for membership management, synchronization,
file compatibility, and usability. Their centralized design
simplifies collaboration and reduces the technical burden on
users. However, these systems depend on centralized
infrastructure for storage and coordination, which introduces
reliance on external services and limits direct peer-to-peer
collaboration.

\subsection{Collaborative Version Control Systems}

Representative examples of collaborative version control
systems include \textit{Git}, \textit{Apache Subversion (SVN)},
and \textit{Perforce}. \textit{Git} is a distributed version
control system in which each participant maintains a local copy
of the repository history. \textit{SVN} follows a centralized
repository model, while \textit{Perforce} is commonly used in
large-scale collaborative projects involving both source code
and binary assets.

\subsubsection{Membership Management}

Collaborative version control systems support collaboration
through shared repositories and controlled contribution
workflows. Users can join projects, obtain repository access,
and contribute changes according to repository permissions.

For example, \textit{Git} repositories can be distributed among
multiple collaborators, while centralized systems such as
\textit{SVN} and \textit{Perforce} typically manage access
through a central repository server.

\subsubsection{Synchronization Model}

Version control systems are not primarily designed for
continuous file synchronization. Instead, changes are usually
exchanged through explicit operations such as commits, pushes,
updates, or pulls.

For example, a modification made in \textit{Git} becomes
available to other collaborators only after the change has
been committed and shared through the repository workflow.
Consequently, synchronization is typically user-driven rather
than automatic.

\subsubsection{History Management}

History management is the primary strength of collaborative
version control systems. These systems maintain detailed
revision histories and provide mechanisms for reviewing,
restoring, and comparing previous versions of files.

For example, \textit{Git} records changes as commits and
supports branching, merging, rollback, and conflict
resolution. Similar capabilities are also provided by
\textit{SVN} and \textit{Perforce} through their respective
revision management workflows.

\subsubsection{File Support}

Version control systems can manage a wide variety of file
types, including source code, documents, images, multimedia
content, datasets, and binary assets.

However, support for large files and frequently changing
binary content varies significantly between systems. For
example, \textit{Git} is primarily optimized for text-based
artifacts such as source code, whereas \textit{Perforce}
provides stronger support for large binary assets commonly
encountered in game development and media production.


\subsubsection{Filesystem Transparency and Usability}

Version control systems typically require users to interact
with repositories through dedicated workflows and tools.
Operations such as committing, updating, merging, and
resolving conflicts introduce additional complexity compared
with ordinary file management.

For example, users of \textit{Git} must explicitly manage
repository state and revision history rather than relying on
fully transparent background synchronization.

\subsection{Peer-to-Peer Synchronization Systems}

Representative examples of peer-to-peer synchronization
systems include \textit{Syncthing} and \textit{Resilio Sync}.
Both systems are designed to synchronize shared folders across
multiple devices without relying on centralized cloud storage.
Instead, participating devices communicate directly with one
another and collectively maintain replicated copies of shared
data.

\subsubsection{Membership Management}

In peer-to-peer synchronization systems, membership management
is usually organized around devices, shared folders, and access
credentials rather than centralized user accounts. A peer must
be authorized before it can participate in the synchronization
of a shared directory.

For example, \textit{Syncthing} manages membership through
device identifiers and folder sharing approvals. A device can
join a synchronized folder only after the relevant peers approve
the sharing relationship. \textit{Resilio Sync} uses secret keys
or sharing links to grant access to synchronized folders, so
possession of the appropriate key determines whether a device
can join the group.

\subsubsection{Synchronization Model}

Synchronization is the primary strength of peer-to-peer
synchronization systems. File changes are automatically
detected and propagated among participating devices without
requiring explicit user intervention.

For example, when a file is modified within a synchronized
folder in \textit{Syncthing}, the update is automatically
distributed to other connected devices. Changes may also be
propagated after temporary disconnections, allowing
participants to continue working while offline and synchronize
once connectivity is restored.

\subsubsection{History Management}

Peer-to-peer synchronization systems generally provide limited
support for history management. Their primary objective is to
maintain consistent copies of files rather than preserve
comprehensive revision histories.

For example, \textit{Syncthing} offers file versioning
mechanisms that can retain previous copies of modified files.
However, these mechanisms are typically based on snapshots or
archived versions rather than structured revision tracking.

As a result, rollback capabilities are usually available, but
conflict resolution and long-term history management remain
less sophisticated than those provided by dedicated version
control systems.

\subsubsection{File Support}

Peer-to-peer synchronization systems support a wide variety of
file types, including source code, documents, images,
multimedia content, datasets, and other binary assets.

For example, both \textit{Syncthing} and \textit{Resilio Sync}
can synchronize heterogeneous project directories containing
mixed file types without requiring specialized handling for
individual formats.

Large files are also supported, making these systems suitable
for collaborative projects involving multimedia content or
other storage-intensive resources.

\subsubsection{Filesystem Transparency and Usability}

Peer-to-peer synchronization systems typically integrate
directly with the local filesystem through ordinary
directories. Users can create, modify, move, and delete files
using their preferred tools while synchronization occurs in
the background.

For example, synchronized folders in \textit{Syncthing}
appear as normal directories on the local machine, allowing
users to work without interacting directly with the
synchronization protocol.

This approach provides strong filesystem transparency and
requires relatively little workflow adaptation compared with
traditional version control systems.

\subsection{Summary of Trade-offs}


\subsection{Evaluation of Existing Paradigms}
\begin{enumerate}
    \item \textbf{Cloud Directory Services} (e.g., Google Drive \cite{googledrive}, Dropbox \cite{dropbox}, Nextcloud): \\
    These services rely on centralized cloud infrastructure. While they offer simple setup and ease of use, they suffer from high transfer latency, storage volume quotas, and strict dependency on external internet connectivity. Centralized hosting also raises user data privacy concerns and imposes access permissions that create collaboration bottlenecks for temporary teams.
    
    \item \textbf{Version Control Systems} (e.g., Git \cite{gitbook, github}, SVN): \\
    VCS tools track repository history using directed acyclic graphs (DAGs). Although they provide powerful conflict resolution, change logs, and branch-level safety, they are designed for text-based source code. They do not handle large binary files well and require users to manually run staging, committing, and remote push/pull cycles. This workflow introduces steep learning curves for non-technical team members.
    
    \item \textbf{P2P File Synchronization} (e.g., Syncthing \cite{syncthing}, Resilio Sync \cite{resilio}): \\
    Decentralized sync programs mirror directories directly between devices without central servers, avoiding cloud storage quotas and allowing high-speed local transfers. However, they lack integrated version rollback histories, making accidental file overwrites difficult to undo. Setup also requires manual node sharing and folder approvals, which adds configuration friction for short-term projects.
\end{enumerate}

\subsection{Key Product Differences: Syncthing vs. Resilio Sync}
We analyze the differences between the leading peer-to-peer synchronization tools to clarify our design choices:
\begin{itemize}
    \item \textbf{Architecture and Openness:} Syncthing is open-source and uses the public Block Exchange Protocol (BEP). It enforces a strict security model where devices must exchange long cryptographic IDs and mutually approve folder invitations. Resilio Sync is proprietary, based on the BitTorrent protocol, and supports selective sync, but has commercial license limits and proprietary lock-in.
    \item \textbf{Lightweight Collaboration Gap:} Both tools act as generic directory mirrors rather than project-specific workspaces. They lack automated, non-destructive version histories that team members can navigate easily. If a peer saves a broken draft, that file immediately overwrites the state of all other online peers, requiring complex administrative configuration to recover.
\end{itemize}

% NOTE: This part should be moved to positioning...
% \subsection{Proposed P2P Synchronization System}
% To address these limitations, we propose a serverless local directory synchronization system. It combines the direct transfer speed of P2P mirroring with the version safety of a lightweight revision tool. The proposed system features:
% \begin{itemize}
%     \item \textbf{Zero Configuration (planned):} Peer discovery relies on Multicast DNS (mDNS) \cite{rfc6762} and DNS-SD \cite{rfc6763} broadcasts on the local subnet, forming a sync mesh without manual IP or node ID exchange.
%     \item \textbf{Local Socket Transfer (planned):} TCP connections \cite{rfc9293} handle data replication directly between peers over the local network, avoiding external server roundtrips.
%     \item \textbf{Autonomous Coordination (planned):} Every node acts as an independent repository coordinator, allowing concurrent edits that propagate without administrative approval bottlenecks.
%     \item \textbf{Durability and Recovery (planned):} We combine a Last-Write-Wins (LWW) conflict policy with automated local backups before any overwrite. This avoids the implementation complexity of vector clocks or CRDTs for small-scale student teams (3–5 peers).
% \end{itemize}

% NOTE: This part should be moved to later section where I don't really know.
% \subsection{Core Project Hypotheses}
% We formulate the following testable scientific hypotheses to guide the evaluation of our proposed system:
% \begin{itemize}
%     \item \textbf{Primary Hypothesis ($H_1$):} A localized, serverless synchronization architecture using mDNS/DNS-SD discovery and direct TCP socket replication can achieve sub-second average synchronization latency on standard local area networks (WLAN) without reliance on external cloud infrastructure.
%     \item \textbf{Secondary Hypothesis ($H_2$):} Utilizing a single logical timestamp (Lamport clock) counter per file, paired with a deterministic Last-Write-Wins (LWW) resolution rule and local file backups, provides sufficient consistency and rollback recovery for small-scale collaborative teams (3-5 peers) while avoiding the state and computational overhead of vector clocks or CRDTs.
% \end{itemize}

% NOTE: This part also should be moved to later section. 
% \subsection{Structured Comparison Table}
% Table \ref{tab:comparison} evaluates the feature tradeoffs of the three paradigms.
%
% \begin{longtable}{>{\raggedright\arraybackslash}p{3.2cm} >{\raggedright\arraybackslash}p{3.8cm} >{\raggedright\arraybackslash}p{3.8cm} >{\raggedright\arraybackslash}p{3.8cm}}
% \caption{Comparison of Collaborative File Sharing and Synchronization Approaches} \label{tab:comparison} \\
% \toprule
% \textbf{Feature / Dimension} & \textbf{Cloud Directory Services} & \textbf{Version Control Systems} & \textbf{P2P File Sync (Proposed)} \\
% \midrule
% \endfirsthead
% \multicolumn{4}{c}{{\bfseries Table \thetable\ -- Continued from previous page}} \\
% \midrule
% \textbf{Feature / Dimension} & \textbf{Cloud Directory Services} & \textbf{Version Control Systems} & \textbf{P2P File Sync (Proposed)} \\
% \midrule
% \endhead
% \midrule
% \multicolumn{4}{r}{{Continued on next page}} \\
% \endfoot
% \bottomrule
% \endlastfoot
% Primary Topology & Centralized Client-Server & Distributed (DAG-based) & Peer-to-Peer (Mesh) \\
% \midrule
% Network Discovery & Centralized cloud portal & Remote repository server & Local mDNS / LAN broadcast \\
% \midrule
% Setup Complexity & Account setup & Repository creation / auth & Low (auto-discovery planned) \\
% \midrule
% Target Use Case & Generic office documents & Source code tracking & Live directory mirroring \\
% \midrule
% Conflict Resolution & Duplicate file generation & Manual merges / diff resolution & LWW (planned) \\
% \midrule
% Version Control & Periodic file revisions & Explicit commits and branches & Local backups (planned) \\
% \midrule
% File Watcher & Periodic polling & Manual staging & Kernel events (planned) \\
% \midrule
% Bandwidth Efficiency & Server upload/download & Incremental compression diffs & Direct sockets / planned delta sync \\
% \midrule
% Internet Dependency & High (cloud servers required) & High (for pull/push) & None (offline-capable) \\
% \midrule
% Data Privacy & Low (third-party storage) & Medium (third-party servers) & High (local-only storage) \\
% \midrule
% Access Permissions & Centralized ACL approvals & Branch rules and PR reviews & Peer-to-peer access \\
% \end{longtable}
